\chapter{Vitamin D (25-Hydroxy Vitamin D)}
\section*{What Is This Test?}
Vitamin D is a fat-soluble secosteroid hormone produced in the skin from 7-dehydrocholesterol upon exposure to ultraviolet B (UVB) radiation (wavelength 290--315 nm) and obtained from the diet (fatty fish, fortified milk) and supplements. The biologically inactive form --- vitamin D3 (cholecalciferol, from skin and animal sources) and vitamin D2 (ergocalciferol, from plant sources and supplements) --- is hydroxylated in the liver by 25-hydroxylase (CYP2R1) to 25-hydroxy vitamin D (25(OH)D, calcidiol), the major circulating form and the best biomarker for vitamin D status. 25(OH)D is further hydroxylated in the kidney by 1-alpha-hydroxylase (CYP27B1) to the biologically active form, 1,25-dihydroxy vitamin D (1,25(OH)2D, calcitriol), which acts through the vitamin D receptor (VDR) to regulate calcium and phosphate homeostasis (increased intestinal absorption, renal reabsorption, and bone resorption), bone mineralization, and immune function. Vitamin D deficiency causes rickets in children (bowed legs, rachitic rosary, craniotabes, delayed fontanelle closure) and osteomalacia in adults (bone pain, proximal muscle weakness, pathologic fractures). Deficiency is common: prevalence of 25(OH)D <20 ng/mL is approximately 30--40\% in the general U.S. population, higher in African Americans (due to melanin blocking UVB), obese individuals (sequestration in adipose tissue), institutionalized elderly, and patients with malabsorption.
\section*{Alternative Names}
25-Hydroxy Vitamin D; 25(OH)D; Calcidiol; Vitamin D3; Cholecalciferol Level; Vitamin D Status
\section*{Principle}
Total 25-hydroxy vitamin D (25(OH)D2 + 25(OH)D3) is measured by LC-MS/MS (gold standard, separately quantifies D2 and D3) or automated competitive chemiluminescent immunoassay (CLIA) on platforms including DiaSorin Liaison, Roche Cobas, Abbott Architect, and Siemens. Immunoassays may over- or under-estimate total 25(OH)D due to differential cross-reactivity with D2, D3, and metabolites. LC-MS/MS is preferred.
\section*{History}
Rickets was described in detail in 17th-century England, where it was known as the ``English disease.'' In 1919--1920, Edward Mellanby produced rickets experimentally in dogs and showed that cod liver oil prevented it, while Huldschinsky demonstrated that ultraviolet light cured rickets in children. In 1922, McCollum and colleagues named the antirachitic fat-soluble factor ``vitamin D,'' distinguishing it from vitamin A. In 1924, Hess and Steenbock independently showed that ultraviolet irradiation endows foods and skin lipids with antirachitic activity, establishing the dual dietary and cutaneous origins of the vitamin. The principal circulating metabolite, 25-hydroxyvitamin D, was identified in 1968 by the laboratory of Hector DeLuca, and the hormonally active form, 1,25-dihydroxyvitamin D, was isolated in 1971 by the groups of DeLuca and Kodicek. The Vitamin D Standardization Program of the 2010s harmonized 25(OH)D assays across laboratories worldwide.
\section*{Why This Test Is Ordered}
\begin{itemize}
  \item Evaluation of suspected vitamin D deficiency in patients with bone pain, proximal muscle weakness, or pathologic fractures
  \item Confirmation of rickets in children and osteomalacia in adults
  \item Assessment of secondary hyperparathyroidism, hypocalcemia, hypophosphatemia, and elevated alkaline phosphatase
  \item Screening of at-risk groups: institutionalized or housebound elderly, individuals with dark skin, obese patients, and patients with malabsorption (celiac disease, Crohn's disease, post-bariatric surgery)
  \item Part of the osteoporosis workup to document status before therapy and to guide supplementation
  \item Monitoring of 25(OH)D levels after 3 months of replacement therapy to confirm adequate repletion
\end{itemize}
\section*{Is This a Primary or Secondary Test}
25(OH)D measurement is the primary diagnostic test for vitamin D status because it reflects total body stores from both cutaneous synthesis and dietary intake. Measurement of 1,25-dihydroxyvitamin D is reserved for specialized disorders of phosphate or calcium metabolism and is not used to assess nutritional status.
\section*{How This Test Is Performed}
A blood sample is collected by standard venipuncture, typically into a plain serum or EDTA tube; no special handling is required and 25(OH)D is stable in serum. In the laboratory, total 25-hydroxyvitamin D (D2 plus D3) is measured by LC-MS/MS or by automated immunoassay. Results are reported in ng/mL or nmol/L, where 1 ng/mL is approximately 2.5 nmol/L.
\section*{What Preparation Is Needed}
No fasting or special preparation is required, and the sample may be drawn at any time of day. The patient should continue their usual diet and medications, and recent sun exposure does not require deferral of testing. If baseline status is being assessed, the test is best performed before starting vitamin D supplementation.
\section*{Contraindications and Precautions}
No contraindications to standard venipuncture. Interpretive cautions: (1) 25(OH)D is the correct test for nutritional status, and 1,25(OH)2D should not be used for this purpose because it is tightly regulated and may be normal or even elevated despite deficiency; (2) assay methods vary, and LC-MS/MS is the reference method, with immunoassays able to over- or underestimate values; (3) severe obesity lowers measured 25(OH)D due to adipose sequestration without necessarily indicating tissue deficiency; and (4) target ranges differ between the Institute of Medicine (adequate $>$20 ng/mL) and Endocrine Society (30--100 ng/mL) guidelines.
\section*{Risks}
Minimal risk. Slight pain or bruising at the venipuncture site; rarely hematoma, infection, or vasovagal syncope.
\section*{What Happens After the Test}
Results are usually available within 1--3 days. Deficient patients are treated with high-dose cholecalciferol (vitamin D3) or ergocalciferol (vitamin D2), typically 50,000 IU weekly for 8 weeks followed by a maintenance dose, and the level is rechecked after 3 months of therapy. Treatment aims for a 25(OH)D level of 30--50 ng/mL, and serum calcium and PTH are followed when secondary hyperparathyroidism is present.
\section*{Understanding Results}
\subsection*{Deficiency (<20 ng/mL, <50 nmol/L per Endocrine Society)}
Rickets (children), osteomalacia (adults), secondary hyperparathyroidism (elevated PTH), hypocalcemia, hypophosphatemia, increased alkaline phosphatase, osteoporosis, increased risk of fractures and falls. Associated with: autoimmune diseases (multiple sclerosis, type 1 diabetes, rheumatoid arthritis), cardiovascular disease, and certain malignancies.
\subsection*{Insufficiency (20--29 ng/mL)}
Increased PTH, bone loss, increased fracture risk. Supplementation recommended.
\subsection*{Sufficiency (30--100 ng/mL, per Endocrine Society)}
Adequate for bone health, suppression of PTH. Optimal calcium absorption.
\subsection*{Toxicity (>100 ng/mL, >150 ng/mL per Endocrine Society)}
Hypercalcemia, hypercalciuria (risk of nephrolithiasis and nephrocalcinosis), vascular calcification, renal failure. Vitamin D toxicity is caused almost exclusively by excessive supplementation (usually >10,000 IU/day for months), NOT by sun exposure.
\section*{Normal Results}
Deficient: <20 ng/mL (<50 nmol/L). Insufficient: 20--29 ng/mL (50--74 nmol/L). Sufficient: 30--100 ng/mL (75--250 nmol/L). Potentially toxic: >100--150 ng/mL. The optimal level for bone health is 30--50 ng/mL per Endocrine Society. The IOM (Institute of Medicine) defines adequacy as >20 ng/mL.
\section*{Follow-up Tests}
Serum calcium, phosphate, PTH (intact), alkaline phosphatase, 24-hour urinary calcium, bone mineral density (DXA), and repeat 25(OH)D after 3 months of supplementation.
\section*{References}
\begin{enumerate}
  \item Holick MF, Binkley NC, Bischoff-Ferrari HA, et al. Evaluation, treatment, and prevention of vitamin D deficiency: an Endocrine Society clinical practice guideline. J Clin Endocrinol Metab. 2011;96(7):1911--1930.
  \item Ross AC, Manson JE, Abrams SA, et al. The 2011 report on dietary reference intakes for calcium and vitamin D from the Institute of Medicine. J Clin Endocrinol Metab. 2011;96(1):53--58.
  \item Holick MF. Vitamin D deficiency. N Engl J Med. 2007;357(3):266--281.
  \item DeLuca HF. Overview of general physiologic features and functions of vitamin D. Am J Clin Nutr. 2004;80(6 Suppl):1689S--1696S.
  \item Dawson-Hughes B, Heaney RP, Holick MF, et al. Estimates of optimal vitamin D status. Osteoporos Int. 2005;16(7):713--716.
\end{enumerate}
\clearpage
